% Unit 072 English translation of chapter6.tex lines 1--320.
% Source: Methods of Algebra, Volume 2, commit
% 9a5803ff77dd3257484cb177f851a73770a59dd3 (CC BY 4.0).
% stable-unit-id: o014.aljabr2.chapter6.overview-g-modules-resolutions
% Coverage: Chapter 6 opening and complete Section 6.1 on G-modules and the
% standard and normalized free resolutions; stops before Section 6.2 at line 321.

% segment-id: o014.li2.chapter6.overview-g-modules-resolutions.h001
\chapter{Group Homology and Cohomology}\label{sec:group-coh}
\hypertarget{o014.aljabr2.chapter6.overview-g-modules-resolutions}{ }

% segment-id: o014.li2.chapter6.overview-g-modules-resolutions.p001
Let $G$ be a group and let $\Bbbk$ be a commutative ring. A $G$-module means
a $\Bbbk$-module equipped with a linear left action of $G$, written as left
multiplication; in many situations $\Bbbk$ is usually taken to be $\Z$. All
$G$-modules form an Abelian category $G\dcate{Mod}$, isomorphic to
$\Bbbk[G]\dcate{Mod}$, where $\Bbbk[G]$ is the group algebra of $G$. For a
$G$-module $M$, define the $\Bbbk$-modules
% segment-id: o014.li2.chapter6.overview-g-modules-resolutions.q001
\begin{align*}
	\text{invariants} \quad M^G & := \left\{x \in M: \forall g \in G, \; gx=x \right\}, \\
	\text{coinvariants} \quad M_G & := M / \lrangle{gx-x: g \in G, x \in M}.
\end{align*}

% segment-id: o014.li2.chapter6.overview-g-modules-resolutions.p002
From the algebraic point of view, the cohomology $\Hm^n(G,M)$ and homology
$\Hm_n(G,M)$ of a $G$-module $M$ are, respectively, the values at $M$ of the
right-derived functors of $(\cdot)^G$ and the left-derived functors of
$(\cdot)_G$, where $n\in\Z_{\geq0}$.

% segment-id: o014.li2.chapter6.overview-g-modules-resolutions.p003
Group homology and cohomology have a long history. Around 1935, W.\ Hurewicz
studied path-connected topological spaces $E$ whose higher homotopy groups
satisfy $n>1\implies\pi_n(E)=0$. He proved that their homology and cohomology
groups are completely determined by the fundamental group $\pi_1(E)$; this is
perhaps the earliest documented formulation of group cohomology. As
homological algebra advanced rapidly during the twentieth century, group
cohomology and homology soon acquired purely algebraic formulations and were
used to address problems within algebra itself. Today they are basic tools in
number theory, topology, and geometry.

% segment-id: o014.li2.chapter6.overview-g-modules-resolutions.p004
On the other hand, $\Hm^n(G,M)$ and $\Hm_n(G,M)$ can also be identified with
$\Ext^n(\Bbbk,M)$ and $\Tor_n(\Bbbk,M)$, respectively, in
$\Bbbk[G]\dcate{Mod}$, with $\Bbbk$ regarded as a trivial $G$-module.
Computing $\Ext$ and $\Tor$ from a projective resolution of $\Bbbk$ gives
another description of $\Hm^n(G,M)$ and $\Hm_n(G,M)$. In fact, we shall give
a canonical resolution $\mathsf{L}\to\Bbbk$ of $\Bbbk$ as a left (or right)
$\Bbbk[G]$-module. It yields the standard cochain complex
$(C^n(G,M))_n$ (or chain complex $(C_n(G,M))_n$) representing $\Hm^n(G,M)$
(or $\Hm_n(G,M)$); see Propositions~\ref{prop:groupcoh-cochain} and
\ref{prop:grouph-chain}. Later, \S\S\ref{sec:homology-computation}---
\ref{sec:bar-resolution} will interpret these canonical resolutions from the
viewpoint of simplicial theory.

% segment-id: o014.li2.chapter6.overview-g-modules-resolutions.p005
The preceding material constitutes the main content of
\S\S\ref{sec:G-mod}---\ref{sec:group-coh-sub}. This chapter adopts two
viewpoints at once: the classical theory of derived functors, especially
universal $\delta$-functors and their characterization, and derived categories.
The classical theory has the advantage of being operational, whereas derived
categories provide finer information than the individual cohomology or
homology groups and often offer a clearer line of thought. The standard
cochain or chain complexes provide the most detailed and canonical
information, since they operate at the level of complexes.

% segment-id: o014.li2.chapter6.overview-g-modules-resolutions.p006
In \S\ref{sec:grp-low-degree}, we shall give concrete interpretations of
$\Hm^1$ and $\Hm^2$: they correspond, respectively, to crossed homomorphisms
and equivalence classes of extensions of $G$ by an Abelian group. The precise
meaning of a group extension is given in Definition~\ref{def:group-ext}. In
fact, if all extensions of $G$ by $M$ are organized into a 2-category
$\cate{Ext}(G,M)$, then the truncation
$\tau^{\leq2}\overline{C}(G,M)$ of the normalized standard complex is exactly
a linear incarnation of $\cate{Ext}(G,M)$; this is the content of
Remark~\ref{rem:Ext-incarnation}.

% segment-id: o014.li2.chapter6.overview-g-modules-resolutions.p007
For any group homomorphism $\varphi:H\to G$ and any $G$-module $M$, its
pullback $\varphi^*M$ can be defined as an $H$-module. General constructions
from module theory give simple descriptions of both the right and left
adjoints of $\varphi^*$. When $\varphi$ is the inclusion of a subgroup,
$\varphi^*$ is called the restriction functor on $G$-modules and is also
written $\Res^G_H$, while its right (respectively left) adjoint is written
$\Ind^G_H$ (respectively $\iInd^G_H$); both are called induction functors.
Their properties are the subject of \S\ref{sec:induced-module}. In particular,
we shall prove the so-called Shapiro lemma (Theorem~\ref{prop:Shapiro}):
% segment-id: o014.li2.chapter6.overview-g-modules-resolutions.q002
\[ \Hm^n\left(G, \Ind^G_H(N)\right) \simeq \Hm^n(H, N), \quad \Hm_n\left(G, \iInd^G_H(N)\right) \simeq \Hm_n(H, N). \]

% segment-id: o014.li2.chapter6.overview-g-modules-resolutions.p008
Continuing the general discussion of $\varphi^*$,
\S\ref{sec:change-of-group} will define families of canonical homomorphisms
% segment-id: o014.li2.chapter6.overview-g-modules-resolutions.q003
\[ \Hm^n(G, M) \to \Hm^n(H, \varphi^* M), \quad \Hm_n(H, \varphi^* M) \to \Hm_n(G, M). \]
These homomorphisms can be combined with the functoriality of cohomology or
homology. Thus every equivariant homomorphism $f:M\to N$ (or $f:N\to M$)
induces the corresponding canonical homomorphism $\Hm^n(f)$ (or $\Hm_n(f)$),
where $M$ is a $G$-module and $N$ an $H$-module; see
Definition~\ref{def:change-of-group-equiv}. In the special case $n=2$,
\S\ref{sec:group-ext-revisited} will interpret these canonical homomorphisms
in terms of group extensions.

% segment-id: o014.li2.chapter6.overview-g-modules-resolutions.p009
For concrete illustrations, \S\ref{sec:coh-cyclic-free} discusses finite
cyclic groups and free groups. The arguments harmoniously combine techniques
from group theory and cohomology. The cohomological and homological dimensions
of groups (Definition~\ref{def:group-dim}) will also enter at the appropriate
point.

% segment-id: o014.li2.chapter6.overview-g-modules-resolutions.p010
The concrete flavor of group theory continues in
\S\ref{sec:group-finite-index}. When $(G:H)$ is finite, we shall state further
properties of the restriction and induction functors; indeed, in this case
$\Ind^G_H=\iInd^G_H$. Definition--Proposition~\ref{def:cor} will give
families of canonical homomorphisms
% segment-id: o014.li2.chapter6.overview-g-modules-resolutions.q004
\[ \mathrm{cor}^n: \Hm^n(H, M) \to \Hm^n(G, M), \quad \mathrm{cor}_n: \Hm_n(G, M) \to \Hm_n(H, M), \]
where $M$ is a $G$-module and $\Res^G_HM$ is abbreviated to $M$ to save
notation. These are called corestriction homomorphisms, and their direction is
opposite to that of the canonical homomorphisms induced by $H\hookrightarrow
G$ (namely restriction, written $\mathrm{res}^n$ and $\mathrm{res}_n$). They
satisfy the fundamental identities
$\mathrm{cor}^n\mathrm{res}^n=(G:H)$ and
$\mathrm{res}_n\mathrm{cor}_n=(G:H)$
(Proposition~\ref{prop:res-cor}), while the special case
$\mathrm{cor}_1:\Hm_1(G,\Z)\to\Hm_1(H,\Z)$ becomes the classical transfer
homomorphism in group theory,
$\mathrm{Ver}_{G|H}:G_{\mathrm{ab}}\to H_{\mathrm{ab}}$.

% segment-id: o014.li2.chapter6.overview-g-modules-resolutions.p011
Let $H$ be a normal subgroup of $G$ and let $M$ be a $G$-module. Then
$\Hm^q(H,M)$ and $\Hm_q(H,M)$ become $(G/H)$-modules. The
Lyndon--Hochschild--Serre spectral sequences discussed in \S\ref{sec:LHS-SS},
% segment-id: o014.li2.chapter6.overview-g-modules-resolutions.q005
\begin{align*}
	E_2^{p, q} = \Hm^p(G/H, \Hm^q(H, M)) & \Rightarrow \Hm^{p+q}(G, M), \\
	E^2_{p, q} = \Hm_p(G/H, \Hm_q(H, M)) & \Rightarrow \Hm_{p+q}(G, M),
\end{align*}
are indispensable advanced tools in the theory of group cohomology and
homology; see Theorem~\ref{prop:LHS-SS}. The information in the exact sequence
of their low-degree terms is extremely useful. The
Lyndon--Hochschild--Serre spectral sequence is at once an application of the
Grothendieck spectral sequence and the result of a canonical filtration on the
standard complex. The latter description gives finer information, such as the
multiplicative structure mentioned in Remark~\ref{rem:LHS-SS-mult}.

% segment-id: o014.li2.chapter6.overview-g-modules-resolutions.p012
More precisely, the multiplicative operation on group cohomology is the cup
product, which originates in topology. In \S\ref{sec:cup-product}, we shall
approach it from two viewpoints. First, the cup product is understood as a
natural reflection of the tensor-product operation on $G$-modules; the
language of derived categories is useful here. Second, its formula is written
directly at the level of the standard complex
(Definition--Proposition~\ref{def:group-cup-3}). The second description is, of
course, an explicit realization of the first. The key to deriving it is to
lift $\Bbbk\rightiso\Bbbk\otimes\Bbbk$ appropriately to a ``diagonal
embedding'' $\Delta:\mathsf{L}\to\mathsf{L}\otimes\mathsf{L}$; several signs
in the argument require careful handling.

% segment-id: o014.li2.chapter6.overview-g-modules-resolutions.p013
If $G$ is finite, the Tate cohomology $\TaHm^n(G,M)$ introduced in
\S\ref{sec:Tate-coh} combines cohomology and homology. Its construction
depends on the canonical homomorphism peculiar to finite groups,
% segment-id: o014.li2.chapter6.overview-g-modules-resolutions.q006
\[ \nu: M_G \to M^G, \quad (\text{the image of }x \in M) \mapsto \sum_{g \in G} gx. \]
For $n\geq1$ (respectively $n\leq-2$), Tate cohomology equals $\Hm^n(G,M)$
(respectively $\Hm_{-n-1}(G,M)$), while for $n=0$ (respectively $n=-1$) it
equals $\Coker(\nu)$ (respectively $\Ker(\nu)$). Tate cohomology continues to
satisfy such properties as long exact sequences and Shapiro's lemma. In many
respects it is even easier to work with than cohomology or homology, and it has
especially prominent applications in number theory.

% segment-id: o014.li2.chapter6.overview-g-modules-resolutions.p014
The preceding theory concerns only discrete groups; groups equipped with a
topology present a much more intricate case. Profinite groups $G$ and smooth
$G$-modules, discussed in \S\ref{sec:profinite-cohomology}, form a special case
that is both simple and useful; the necessary background is given in
\S\ref{sec:profinite-groups}. The theory remains essentially algebraic. The
corresponding cohomology can be understood as a right-derived functor or,
equivalently, defined as a $\varinjlim$ of the cohomology of finite groups; the
continuous version of the standard complex provides another method of
calculation. In number theory, the main cases of interest are those in which
$G$ is a Galois group or an arithmetic fundamental group.

% segment-id: o014.li2.chapter6.overview-g-modules-resolutions.p015
Finally, \S\ref{sec:nonabelian-coh} introduces non-Abelian group cohomology.
For a non-Abelian group, usually only $\Hm^0$ and $\Hm^1$ can be defined, so
its long exact sequence is correspondingly limited
(Theorem~\ref{prop:nonabelian-long}). This theory is particularly well suited
to classifying geometric or algebraic objects, because the symmetry groups of
such objects are often non-Abelian. Several examples will be discussed later
in \S\ref{sec:Galois-H1}. These constructions also apply to profinite groups.

% segment-id: o014.li2.chapter6.overview-g-modules-resolutions.n001
\begin{petunjukbacaan}
	The role of the coefficient ring $\Bbbk$ in this chapter is secondary, and
	it is often taken to be $\Z$; see the explanation in
	Remark~\ref{rem:Gmod-k}. Readers may adapt the material to their needs and
	skip arguments involving derived categories; doing so will not interfere
	with most classical applications. In particular, the cup product on
	cohomology may be treated entirely through the explicit operations in
	Definition--Proposition~\ref{def:group-cup-3}. The notion of a $G$-module
	(see \S\ref{sec:G-mod}) occasionally appears as an example or exercise in
	later chapters and usually involves only the basic definitions. The
	exceptions are \S\ref{sec:descent-Galois} and \S\ref{sec:Galois-H1}, which
	use smooth $G$-modules for profinite groups and their cohomology. Readers
	interested in number theory should in addition master
	\S\S\ref{sec:Tate-coh}---\ref{sec:nonabelian-coh}.
\end{petunjukbacaan}

% segment-id: o014.li2.chapter6.overview-g-modules-resolutions.h002
\section{\texorpdfstring{$G$}{G}-Modules and Their Resolutions}\label{sec:G-mod}

% segment-id: o014.li2.chapter6.overview-g-modules-resolutions.p016
Unless otherwise stated, throughout this section $\Bbbk$ is a commutative ring
and $G$ a group. Write the identity element of $G$ as $1_G$.

% segment-id: o014.li2.chapter6.overview-g-modules-resolutions.d001
\begin{definition}\label{def:G-mod}
	\index{$G$-module}
	\index[sym1]{G-Mod@$G\dcate{Mod}$}
	A \emph{$G$-module} with coefficients in $\Bbbk$ is a $\Bbbk$-module $M$
	equipped with a left action of $G$. The action $G\times M\to M$ is written
	as multiplication and satisfies
% segment-id: o014.li2.chapter6.overview-g-modules-resolutions.q007
	\begin{align*}
		g(m + m') & = gm + gm', \\
		g(tm) & = t(gm),
	\end{align*}
	for all $g\in G$, $t\in\Bbbk$, and $m,m'\in M$. In other words, $M$ is
	equipped with a linear action of $G$.

	A homomorphism $f:M_1\to M_2$ between two $G$-modules is defined to be a
	$\Bbbk$-module homomorphism satisfying the equivariance property
	$f(gm_1)=gf(m_1)$ for all $g\in G$ and $m_1\in M_1$.
\end{definition}

% segment-id: o014.li2.chapter6.overview-g-modules-resolutions.p017
Strictly speaking, the structure above should be called a left $G$-module. A
right $G$-module is defined in an entirely similar way and is equivalent to a
left $G^{\opp}$-module. When there is no danger of confusion, the coefficient
ring $\Bbbk$ is suppressed from the notation, and the category of all
$G$-modules is written $G\dcate{Mod}$. This category is $\Bbbk$-linear
(Definition~\ref{def:k-linear-cat}).

% segment-id: o014.li2.chapter6.overview-g-modules-resolutions.ex001
\begin{example}\label{eg:trivial-G-mod}
	\index{$G$-module!trivial}
	Equip $\Bbbk$ with the trivial action of $G$, namely $gm=m$. The resulting
	object is called the \emph{trivial $G$-module} and is denoted simply by
	$\Bbbk$.
\end{example}

% segment-id: o014.li2.chapter6.overview-g-modules-resolutions.d002
\begin{definition}
	Let $M$ be a $G$-module. If a $\Bbbk$-submodule $N\subset M$ is closed
	under the action of $G$, then $N$ is called a $G$-submodule. Equip the
	quotient $\Bbbk$-module $M/N$ with the action $g(m+N)=gm+N$; this is called
	the corresponding quotient $G$-module.
\end{definition}

% segment-id: o014.li2.chapter6.overview-g-modules-resolutions.p018
The group algebra $\Bbbk[G]$ can be constructed from $G$; see
\cite[Definition 5.6.3]{Li1}. Each of its elements has a unique expression of
the form $\sum_{g\in G}a_gg$, with only finitely many nonzero coefficients
$a_g\in\Bbbk$. From now on, we regard $G$ as a subset of $\Bbbk[G]$.

% segment-id: o014.li2.chapter6.overview-g-modules-resolutions.pr001
\begin{proposition}\label{prop:G-mod-kG}
	There is an isomorphism of categories
	$G\dcate{Mod}\rightiso\Bbbk[G]\dcate{Mod}$ as follows. It preserves the
	underlying $\Bbbk$-module. If $M$ is a $G$-module, its scalar
	multiplication as a left $\Bbbk[G]$-module comes from the homomorphism
% segment-id: o014.li2.chapter6.overview-g-modules-resolutions.q008
	\[ \Bbbk[G] \dotimes{\Bbbk} M \to M, \quad (\sum_g a_g g) \otimes m \mapsto \sum_g a_g (gm), \]
	while $G$-module homomorphisms correspond to $\Bbbk[G]$-module
	homomorphisms.
\end{proposition}

% segment-id: o014.li2.chapter6.overview-g-modules-resolutions.pf001
\begin{proof}
	The inverse functor is determined as follows. For a left $\Bbbk[G]$-module
	$M$, use the embedding of multiplicative monoids
	$G\hookrightarrow\Bbbk[G]$ to let $G$ act on $M$ from the left, thereby
	making $M$ a $G$-module. It is straightforward to verify that both
	directions are well defined and inverse to each other.
\end{proof}

% segment-id: o014.li2.chapter6.overview-g-modules-resolutions.p019
Consequently, $G\dcate{Mod}$ is a $\Bbbk$-linear Abelian category, since
$\Bbbk[G]\dcate{Mod}$ is one. Properties of $G$-submodules and quotient
$G$-modules translate directly into properties of $\Bbbk[G]$-modules.

% segment-id: o014.li2.chapter6.overview-g-modules-resolutions.d003
\begin{definition}\label{def:HomG}
	\index[sym1]{HomG@$\Hom_G$}
	Write $\Hom$ in $G\dcate{Mod}$ as $\Hom_G$. Under the correspondence of
	Proposition~\ref{prop:G-mod-kG}, this is also $\Hom_{\Bbbk[G]}$ in
	$\Bbbk[G]\dcate{Mod}$.
\end{definition}

% segment-id: o014.li2.chapter6.overview-g-modules-resolutions.p020
Several basic operations on $\Bbbk$-modules extend readily to $G$-modules:
% segment-id: o014.li2.chapter6.overview-g-modules-resolutions.l001
\begin{description}
% segment-id: o014.li2.chapter6.overview-g-modules-resolutions.i001
	\item[Direct sums and direct products] For a family of $G$-modules
	$(M_i)_{i\in I}$, equip the $\Bbbk$-modules $\prod_{i\in I}M_i$ and
	$\bigoplus_{i\in I}M_i$ with the left $G$-action
% segment-id: o014.li2.chapter6.overview-g-modules-resolutions.q009
	\[ g (m_i)_{i \in I} := (gm_i)_{i \in I}, \]
	called the diagonal action. With this action, $\prod_{i\in I}M_i$ and
	$\bigoplus_{i\in I}M_i$ are $G$-modules.
% segment-id: o014.li2.chapter6.overview-g-modules-resolutions.i002
	\item[Tensor products] For $G$-modules $M$ and $N$, equip
	$M\otimes N:=M\dotimes{\Bbbk}N$ with the diagonal left $G$-action
% segment-id: o014.li2.chapter6.overview-g-modules-resolutions.q010
	\[ g (x \otimes y) = gx \otimes gy, \quad x \in M, \quad y \in N ; \]
	since $(x,y)\mapsto gx\otimes gy$ is bilinear, this formula is well
	defined. Plainly $x\otimes y\mapsto y\otimes x$ defines an isomorphism of
	$G$-modules $M\otimes N\simeq N\otimes M$.
% segment-id: o014.li2.chapter6.overview-g-modules-resolutions.i003
	\item[$\Hom$ modules] For $G$-modules $M$ and $N$, the $\Bbbk$-module
	homomorphisms between them form the $\Bbbk$-module
	$\Hom(M,N):=\Hom_{\Bbbk}(M,N)$. Equip $\Hom(M,N)$ with the left
	$G$-action
% segment-id: o014.li2.chapter6.overview-g-modules-resolutions.q011
	\[ ({}^g f)(x) = g(f(g^{-1} x)), \quad f \in \Hom(M, N), \quad x \in M, \]
	where $g^{-1}$ and $g$ on the right arise from the left $G$-actions on $M$
	and $N$, respectively. As a composite of three $\Bbbk$-module
	homomorphisms, ${}^gf$ belongs to $\Hom(M,N)$. Thus ${}^gf$ can be viewed
	as the ``conjugate'' $gfg^{-1}$, making the properties required for a
	$G$-module structure evident.

% segment-id: o014.li2.chapter6.overview-g-modules-resolutions.i004
	\item[Contragredient modules] In the construction of the $\Hom$ module,
	take $N=\Bbbk$ with its trivial $G$-action. Then $\Hom(M,\Bbbk)$ becomes a
	$G$-module, called the contragredient module of $M$. Its $G$-action is
	given by ${}^gf=f\circ g^{-1}$.
	\index{$G$-module!contragredient}
\end{description}

% segment-id: o014.li2.chapter6.overview-g-modules-resolutions.p021
All these constructions are functorial in $M$ and $N$. The definitions
immediately give
% segment-id: o014.li2.chapter6.overview-g-modules-resolutions.q012
\begin{align*}
	\Hom_G(M, N) & = \Hom(M, N)^G \\
	& := \left\{f \in \Hom(M, N): \forall g \in G, \; {}^g f = f \right\}.
\end{align*}

% segment-id: o014.li2.chapter6.overview-g-modules-resolutions.p022
Under the isomorphism of categories in Proposition~\ref{prop:G-mod-kG}, direct
sums and direct products of $G$-modules are, respectively, the direct sums and
direct products of $\Bbbk[G]$-modules. The $G$-module structures on tensor
products and $\Hom$ modules are slightly subtler; from the viewpoint of module
theory, they in fact correspond to the Hopf-algebra structure on $\Bbbk[G]$
(Example~\ref{eg:group-Hopf-algebra}). Interested readers may compare the
related discussions in \S\ref{sec:bialgebra} and
\S\ref{sec:duality-examples}; we shall not elaborate here.

% segment-id: o014.li2.chapter6.overview-g-modules-resolutions.pr002
\begin{proposition}\label{prop:Gmod-monoidal}
	With respect to the tensor-product operation, $G\dcate{Mod}$ is a symmetric
	monoidal category whose unit object is the trivial $G$-module $\Bbbk$.
\end{proposition}

% segment-id: o014.li2.chapter6.overview-g-modules-resolutions.pf002
\begin{proof}
	This follows directly from the definitions above.
\end{proof}

% segment-id: o014.li2.chapter6.overview-g-modules-resolutions.p023
Moreover, the usual canonical isomorphism in $\Bbbk\dcate{Mod}$,
% segment-id: o014.li2.chapter6.overview-g-modules-resolutions.g001
\begin{equation*}
	\begin{tikzcd}[row sep=tiny]
		\Hom(L \otimes M, N) \arrow[leftrightarrow, r, "\sim"] & \Hom(L, \Hom(M, N)) \\
		\varphi \arrow[mapsto, r] & {\left[ w \mapsto \varphi(w \otimes (\cdot)) \right]} \\
		{\left[ w \otimes x \mapsto \psi(w)(x) \right]} & \psi \arrow[mapsto, l]
	\end{tikzcd}
\end{equation*}
also lifts to the level of $G$-modules. If $L$, $M$, and $N$ are $G$-modules
and the tensor product and $\Hom$ above carry the $G$-module structures just
defined, then both directions are isomorphisms of $G$-modules; readers are
invited to verify this carefully.

% segment-id: o014.li2.chapter6.overview-g-modules-resolutions.p024
Consequently, ${}^g\varphi=\varphi$ is equivalent to ${}^g\psi=\psi$, giving
the adjunction
% segment-id: o014.li2.chapter6.overview-g-modules-resolutions.e001
\begin{equation}\label{eqn:G-mod-tensor-Hom}
	\begin{tikzcd}
		\Hom_G(L \otimes M, N) \arrow[leftrightarrow, r, "\sim"] & \Hom_G(L, \Hom(M, N)),
	\end{tikzcd}
\end{equation}
where $L$, $M$, and $N$ are all $G$-modules. Here is a simple application.

% segment-id: o014.li2.chapter6.overview-g-modules-resolutions.pr003
\begin{proposition}\label{prop:group-tensor-proj}
	Let $L$ and $M$ be $G$-modules, with $L$ projective as a $G$-module and
	$M$ projective as a $\Bbbk$-module. Then $L\otimes M$ is projective as a
	$G$-module.
\end{proposition}

% segment-id: o014.li2.chapter6.overview-g-modules-resolutions.pf003
\begin{proof}
	The adjunction \eqref{eqn:G-mod-tensor-Hom} shows that the functor
	$(\cdot)\otimes M:G\dcate{Mod}\to G\dcate{Mod}$ has the exact right
	adjoint $\Hom(M,\cdot)$. Proposition~\ref{prop:adjoint-injective-projective}
	therefore implies that $(\cdot)\otimes M$ preserves projective objects.
\end{proof}

% segment-id: o014.li2.chapter6.overview-g-modules-resolutions.d004
\begin{definition}[Restriction and inflation]\label{def:Res-Infl}
	\index{$G$-module!restriction/inflation}
	\index[sym1]{ResGH@$\Res^G_H$}
	\index[sym1]{InflGH@$\mathrm{Infl}^G_{G/H}$}
	For every group homomorphism $\varphi:H\to G$, there is a corresponding
	functor
% segment-id: o014.li2.chapter6.overview-g-modules-resolutions.q013
	\[ \varphi^*: G\dcate{Mod} \to H\dcate{Mod}; \]
	for a $G$-module $M$, its image under $\varphi^*$ is the same
	$\Bbbk$-module $M$, with the $H$-action given by $hm:=\varphi(h)m$. Two
	important cases of the functor $\varphi^*$ are:
% segment-id: o014.li2.chapter6.overview-g-modules-resolutions.l002
	\begin{description}
% segment-id: o014.li2.chapter6.overview-g-modules-resolutions.i005
		\item[Restriction] For a subgroup $H\subset G$, take
		$\varphi:H\hookrightarrow G$. The corresponding functor
		$\Res^G_H:G\dcate{Mod}\to H\dcate{Mod}$ restricts the group action to
		$H$.
% segment-id: o014.li2.chapter6.overview-g-modules-resolutions.i006
		\item[Inflation] For a normal subgroup $H\lhd G$, take
		$\varphi:G\twoheadrightarrow G/H$. The corresponding functor
		$\mathrm{Infl}^G_{G/H}:(G/H)\dcate{Mod}\to G\dcate{Mod}$ pulls the
		group action back to $G$.
	\end{description}
\end{definition}

% segment-id: o014.li2.chapter6.overview-g-modules-resolutions.p025
Notice that if there is a sequence of group homomorphisms
$H\xrightarrow{\varphi}G\xrightarrow{\psi}F$, then the corresponding functors
satisfy $\varphi^*\psi^*=(\psi\varphi)^*$.

% segment-id: o014.li2.chapter6.overview-g-modules-resolutions.p026
From the group-algebra viewpoint, the homomorphism $\varphi:H\to G$ extends to
a $\Bbbk$-algebra homomorphism $\Bbbk[H]\to\Bbbk[G]$, and the corresponding
functor $\Bbbk[G]\dcate{Mod}\to\Bbbk[H]\dcate{Mod}$ is precisely
$\varphi^*$. This functor can also be expressed by a tensor product as
% segment-id: o014.li2.chapter6.overview-g-modules-resolutions.e002
\begin{equation}\label{eqn:G-module-pullback-tensor}
	\varphi^* M \simeq \Bbbk[G] \dotimes{\Bbbk[G]} M, \quad \Bbbk[G] \;\text{regarded via $\varphi$ as a bimodule}\; (\Bbbk[H], \Bbbk[G]).
\end{equation}

% segment-id: o014.li2.chapter6.overview-g-modules-resolutions.p027
The left and right adjoints of $\varphi^*$, together with the corresponding
adjunction isomorphisms, can be written explicitly.

% segment-id: o014.li2.chapter6.overview-g-modules-resolutions.pr004
\begin{proposition}\label{prop:pullback-two-adjoint}
	Use the notation above. For a group homomorphism $\varphi:H\to G$, any
	$G$-module $M$, and any $H$-module $N$, there are canonical isomorphisms of
	$\Bbbk$-modules
% segment-id: o014.li2.chapter6.overview-g-modules-resolutions.g002
	\[\begin{tikzcd}[row sep=tiny]
		\Hom_H(\varphi^* M, N) \arrow[leftrightarrow, r, "\sim"] & \Hom_G\left( M, \Hom_H(\Bbbk[G], N) \right) \\
		\theta \arrow[mapsto, r] & {[x \mapsto [G \ni g \mapsto \theta(gx)] ]} \\
		{[x \mapsto \xi(x)(1_G)]} & \xi, \arrow[mapsto, l] \\
		\Hom_H(N, \varphi^* M) \arrow[leftrightarrow, r, "\sim"] & \Hom_G\left( \Bbbk[G] \dotimes{\Bbbk[H]} N, M \right) \\
		\psi \arrow[mapsto, r] & {[g \otimes y \mapsto g\psi(y)]} \\
		{[y \mapsto \eta(1_G \otimes y) ]} & \eta \arrow[mapsto, l].
	 \end{tikzcd}\]
	Here $\Bbbk[G]$ is regarded both as a
	$(\Bbbk[H],\Bbbk[G])$-bimodule and as a
	$(\Bbbk[G],\Bbbk[H])$-bimodule. In this way,
	$\Hom_H(\Bbbk[G],N)$ and $\Bbbk[G]\dotimes{\Bbbk[H]}N$ acquire
	$G$-module structures.
\end{proposition}

% segment-id: o014.li2.chapter6.overview-g-modules-resolutions.pf004
\begin{proof}
	Expand the definitions directly to verify that both pairs of maps are
	well-defined $\Bbbk$-module homomorphisms; the calculation is not
	difficult. Once this has been checked, it is clear that the maps are
	mutually inverse. Alternatively, apply the general theory in
	\cite[Corollary 6.6.8]{Li1}.
\end{proof}

% segment-id: o014.li2.chapter6.overview-g-modules-resolutions.p028
Since $G\dcate{Mod}$ is an Abelian category, projective and injective
resolutions of $G$-modules can be considered in the usual way. For the subject
of this chapter, the most important are projective resolutions of the trivial
$G$-module $\Bbbk$. Such resolutions not only exist; they also admit a
canonical, explicit choice.

% segment-id: o014.li2.chapter6.overview-g-modules-resolutions.d005
\begin{definition}\label{def:resolution-trivial-module}
	Let $G$ be a group. For each $n\in\Z_{\geq0}$, define
% segment-id: o014.li2.chapter6.overview-g-modules-resolutions.q014
	\[ \mathsf{L}_n := \text{the free $\Bbbk$-module with basis $G^{n+1}$}, \]
	whose elements are written as $\Bbbk$-linear combinations of
	$(g_0,\ldots,g_n)\in G^{n+1}$. For every $n\geq1$, define homomorphisms
% segment-id: o014.li2.chapter6.overview-g-modules-resolutions.q015
	\begin{align*}
		\partial_n, \partial'_n: \mathsf{L}_n & \to \mathsf{L}_{n-1}, \\
		\partial_n\left( g_0, \ldots, g_n \right) & = \sum_{k=0}^n (-1)^{n-k} (\ldots, \widehat{g_k}, \ldots), \\
		\partial'_n\left( g_0, \ldots, g_n \right) & = \sum_{k=0}^n (-1)^k (\ldots, \widehat{g_k}, \ldots),
	\end{align*}
	where $\widehat{g_k}$ indicates that the term is omitted; thus
	$\partial'_n=(-1)^n\partial_n$. Also define the homomorphism
% segment-id: o014.li2.chapter6.overview-g-modules-resolutions.q016
	\[ \partial_0 = \partial'_0: \mathsf{L}_0 \to \Bbbk, \quad (g_0) \mapsto 1. \]

	Define the action
	$g(g_0,\ldots,g_n)=(gg_0,\ldots,gg_n)$ (respectively
	$(g_0,\ldots,g_n)g=(g_0g,\ldots,g_ng)$) so that $\mathsf{L}_n$ becomes a
	left (respectively right) $G$-module, and equip $\Bbbk$ with the trivial
	$G$-action. Then every $\partial_n$ and $\partial'_n$ is a homomorphism of
	left (respectively right) $G$-modules.
\end{definition}

% segment-id: o014.li2.chapter6.overview-g-modules-resolutions.p029
Plainly $(\mathsf{L}_n,\partial_n)_n$ and
$(\mathsf{L}_n,\partial'_n)_n$ differ only by the reversal
$(g_0,\ldots,g_n)\mapsto(g_n,\ldots,g_0)$. This reversal commutes with both the
left and the right $G$-actions, so for all the properties below it suffices to
treat either $\partial_n$ or $\partial'_n$.

% segment-id: o014.li2.chapter6.overview-g-modules-resolutions.r001
\begin{remark}\label{rem:augmentation-kG}
	\index{augmentation homomorphism}
	If $\mathsf{L}_0$ is identified with $\Bbbk[G]$, then
	$\partial_0=\partial'_0:\mathsf{L}_0\to\Bbbk$ is identified with the map
% segment-id: o014.li2.chapter6.overview-g-modules-resolutions.q017
	\[ \Bbbk[G] \to \Bbbk, \quad \sum_g a_g g \mapsto \sum_g a_g; \]
	this is plainly also a surjective $\Bbbk$-algebra homomorphism. It is called
	the \emph{augmentation homomorphism} of $\Bbbk[G]$, and its kernel
% segment-id: o014.li2.chapter6.overview-g-modules-resolutions.q018
	\[ \mathfrak{I} := \left\{ \sum_g a_g g : \sum_g a_g = 0 \right\} \]
	is called the \emph{augmentation ideal} of $\Bbbk[G]$. As a $\Bbbk$-module,
	$\mathfrak{I}$ plainly decomposes with respect to the following basis:
% segment-id: o014.li2.chapter6.overview-g-modules-resolutions.q019
	\[ \mathfrak{I} = \bigoplus_{\substack{g \in G \\ g \neq 1_G}} \Bbbk (g - 1_G). \]
	\index{augmentation ideal}
	\index[sym1]{I-aug@$\mathfrak{I}$}
\end{remark}

% segment-id: o014.li2.chapter6.overview-g-modules-resolutions.le001
\begin{lemma}\label{prop:L-acyclic}
	The complex
	$\cdots\to\mathsf{L}_n\xrightarrow{\partial_n}\cdots
	\xrightarrow{\partial_0}\Bbbk\to0$ defined above is exact. Regarded as a
	complex of $\Bbbk$-modules, its identity morphism is null-homotopic. The
	same assertions hold when $\partial_n$ is replaced by $\partial'_n$.
\end{lemma}

% segment-id: o014.li2.chapter6.overview-g-modules-resolutions.pf005
\begin{proof}
	It suffices to discuss the $\partial'_n$ version. For every $n\geq1$,
	observe that $\partial'_{n-1}\partial'_n$ sends $(g_0,\ldots,g_n)$ to a
	sum of elements of the form
% segment-id: o014.li2.chapter6.overview-g-modules-resolutions.q020
	\[ \pm (\ldots, \widehat{g_p}, \ldots, \widehat{g_q}, \ldots), \quad 0 \leq p < q \leq n. \]
	The signs cancel in pairs, giving a complex; the elementary verification is
	left to the reader.

	Next, for every $n\geq0$, define a $\Bbbk$-module homomorphism
% segment-id: o014.li2.chapter6.overview-g-modules-resolutions.q021
	\[ s_n: \mathsf{L}_n \to \mathsf{L}_{n+1}, \quad (g_0, \ldots, g_n) \mapsto (1_G, g_0, \ldots, g_n), \]
	and define $s_{-1}:\Bbbk\to\mathsf{L}_0$ by sending $1\in\Bbbk$ to
	$1_G\in\mathsf{L}_0$. It is easy to see that
% segment-id: o014.li2.chapter6.overview-g-modules-resolutions.q022
	\[ \partial'_{n+1} s_n + s_{n-1} \partial'_n = \identity_{\mathsf{L}_n} \]
	for every $n\geq-1$, with the conventions $\mathsf{L}_{-1}:=\Bbbk$ and
	$\partial'_{-1}:=0$. This shows that
% segment-id: o014.li2.chapter6.overview-g-modules-resolutions.q023
	\[ \cdots \to \mathsf{L}_n \xrightarrow{\partial'_n} \cdots \xrightarrow{\partial'_0} \Bbbk \to 0 \]
	has a null-homotopic identity morphism as a complex of $\Bbbk$-modules and
	is therefore exact.
\end{proof}

% segment-id: o014.li2.chapter6.overview-g-modules-resolutions.x001
\index[sym1]{L-cplx@$\mathsf{L}$, $\overline{\mathsf{L}}$}

% segment-id: o014.li2.chapter6.overview-g-modules-resolutions.p030
Write the complex
$\cdots\to\mathsf{L}_n\xrightarrow{\partial_n}\cdots\to\mathsf{L}_0\to0$
as $\mathsf{L}$. It may be viewed either as a complex in degrees
$\ldots,-n,\ldots,0$ or as a chain complex in degrees
$\ldots,n,\ldots,0$. In either interpretation, $\partial_0=\partial'_0$
gives a morphism $\mathsf{L}\to\Bbbk$.

% segment-id: o014.li2.chapter6.overview-g-modules-resolutions.p031
This complex and morphism can be understood either as a complex of left
$G$-modules or as a complex of right $G$-modules. By symmetry, we shall mainly
state the left-handed version.

% segment-id: o014.li2.chapter6.overview-g-modules-resolutions.pr005
\begin{proposition}\label{prop:trivial-mod-resolution}
	The morphism $\mathsf{L}\to\Bbbk$ defined above gives a free resolution of
	the trivial $G$-module $\Bbbk$.
\end{proposition}

% segment-id: o014.li2.chapter6.overview-g-modules-resolutions.pf006
\begin{proof}
	As $(h_1,\ldots,h_n)\in G^n$ varies,
	$(1_G,h_1,\ldots,h_n)$ spans a basis of $\mathsf{L}^n$ as a left
	$\Bbbk[G]$-module.
\end{proof}

% segment-id: o014.li2.chapter6.overview-g-modules-resolutions.p032
There is also a canonical free resolution of $\Bbbk$ that is more concise and
sometimes more practical than $\mathsf{L}$.

% segment-id: o014.li2.chapter6.overview-g-modules-resolutions.dp001
\begin{definition-proposition}[Normalized complex]
	For every $n>0$, define $\mathrm{Triv}_n\subset\mathsf{L}_n$ to be the
	$G$-submodule generated by the elements
% segment-id: o014.li2.chapter6.overview-g-modules-resolutions.q024
	\[ (g_0, \ldots, g_n) \in G^{n+1}, \quad \exists\, 0 \leq k < n, \; g_k = g_{k+1}. \]
	Also define $\mathrm{Triv}_0=0$. With respect to either $\partial_n$ or
	$\partial'_n$, these objects form a subcomplex $\mathrm{Triv}$ of
	$\mathsf{L}$. The normalization of $\mathsf{L}$ is then defined to be the
	quotient complex
	$\overline{\mathsf{L}}:=\mathsf{L}/\mathrm{Triv}$. Regarded as a complex of
	$\Bbbk$-modules, the identity morphism on $\mathrm{Triv}$ is
	null-homotopic.
\end{definition-proposition}

% segment-id: o014.li2.chapter6.overview-g-modules-resolutions.pf007
\begin{proof}
	Plainly $\mathrm{Triv}_n$ is invariant under
	$(g_0,\ldots,g_n)\mapsto(g_n,\ldots,g_0)$, so the $\partial_n$ and
	$\partial'_n$ versions do not differ; it suffices to treat the latter. If
	$(g_0,\ldots,g_n)$ satisfies $g_k=g_{k+1}$, then in
	$\partial'_n(g_0,\ldots,g_n)=
	\sum_i(-1)^i(\ldots,\widehat{g_i},\ldots)$, the terms corresponding to
	$i=k,k+1$ cancel, while the remaining terms lie in
	$\mathrm{Triv}_{n-1}$. This proves that the objects form a subcomplex.

	To show that the identity morphism of $\mathrm{Triv}$ is null-homotopic as
	a complex of $\Bbbk$-modules, simply restrict the homotopy
	$s_n:\mathsf{L}_n\to\mathsf{L}_{n+1}$ in the proof of
	Lemma~\ref{prop:L-acyclic} to $\mathrm{Triv}_n$, for $n\geq0$.
\end{proof}

% segment-id: o014.li2.chapter6.overview-g-modules-resolutions.p033
By construction, $\mathsf{L}\to\Bbbk$ vanishes on $\mathrm{Triv}$ and hence
factors through $\overline{\mathsf{L}}\to\Bbbk$.

% segment-id: o014.li2.chapter6.overview-g-modules-resolutions.pr006
\begin{proposition}\label{prop:trivial-mod-nor-resolution}
	The morphism $\overline{\mathsf{L}}\to\Bbbk$ defined above gives a free
	resolution of the trivial $G$-module $\Bbbk$.
\end{proposition}

% segment-id: o014.li2.chapter6.overview-g-modules-resolutions.pf008
\begin{proof}
	First, every
	$\overline{\mathsf{L}}_n=\mathsf{L}_n/\mathrm{Triv}_n$ is a free
	$G$-module: a basis is given by the images of
	$(1_G,h_1,\ldots,h_n)\in G^{n+1}$ with $h_1\neq1_G$ and
	$h_k\neq h_{k+1}$ for $1\leq k<n$.

	Second, the acyclicity of $\mathrm{Triv}$ implies that
	$\mathsf{L}\to\overline{\mathsf{L}}$ is a quasi-isomorphism. Hence
	Proposition~\ref{prop:trivial-mod-resolution} implies that
	$\overline{\mathsf{L}}\to\Bbbk$ is a quasi-isomorphism.
\end{proof}

% segment-id: o014.li2.chapter6.overview-g-modules-resolutions.p034
For later applications, it is more convenient to express the basis in the
following ``bar'' notation. First, for integers $a\leq b$ and a sequence of
elements $(g_a,\ldots,g_b)\in G^{b-a+1}$, introduce the abbreviation
% segment-id: o014.li2.chapter6.overview-g-modules-resolutions.q025
\[ g_{[a, b]} := g_a g_{a+1} \cdots g_b. \]
As a left $G$-module, a basis of $\mathsf{L}_n$ may be taken to be
% segment-id: o014.li2.chapter6.overview-g-modules-resolutions.e003
\begin{equation}\label{eqn:L-basis}
	(g_1 | \cdots | g_n) := \left( 1_G, g_{[1, 1]}, g_{[1, 2]}, \ldots, g_{[1, n]} \right), \quad (g_1, \ldots, g_n) \in G^n.
\end{equation}

% segment-id: o014.li2.chapter6.overview-g-modules-resolutions.p035
Observe that
% segment-id: o014.li2.chapter6.overview-g-modules-resolutions.e004
\begin{equation}\label{eqn:L-basis-partial}\begin{aligned}
	\partial_n (g_1 | \ldots | g_n) & = (-1)^n \left( g_{[1, 1]}, \ldots, g_{[1, n]}\right) + (-1)^{n-1} \left(1_G, g_{[1, 2]}, \ldots, g_{[1, n]}\right) \\
	& \quad + \cdots + \left( 1_G, g_{[1, 1]}, \ldots, g_{[1, n-1]}\right) \\
	& = (-1)^n g_1 (g_2 | \cdots | g_n) + \sum_{k=1}^{n-1} (-1)^{n-k} (\cdots | g_k g_{k+1}| \cdots) \\
	& \quad + (g_1 | \cdots |g_{n-1}), \\
	\partial'_n(g_1 | \cdots | g_n) & = g_1 (g_2 | \cdots | g_n) + \sum_{k=1}^{n-1} (-1)^k (\cdots | g_k g_{k+1}| \cdots) \\
	& \quad + (-1)^n (g_1 | \cdots |g_{n-1}).
\end{aligned}\end{equation}

% segment-id: o014.li2.chapter6.overview-g-modules-resolutions.p036
Now consider the mirror image of \eqref{eqn:L-basis}. Take as a basis of
$\mathsf{L}_n$ as a right $G$-module
% segment-id: o014.li2.chapter6.overview-g-modules-resolutions.e005
\begin{equation}\label{eqn:L-basis-right}
	(g_1 | \cdots | g_n)^\dagger := \left( g_{[1, n]}, \ldots, g_{[n-1, n]} , g_{[n, n]}, 1_G \right), \quad (g_1, \ldots, g_n) \in G^n.
\end{equation}
Then
$\partial'_n(g_1|\cdots|g_n)^\dagger=
(-1)^n\partial_n(g_1|\cdots|g_n)^\dagger$ equals
% segment-id: o014.li2.chapter6.overview-g-modules-resolutions.e006
\begin{multline}\label{eqn:L-basis-partial-right}
	\left( g_{[2, n]}, \ldots, g_{[n, n]}, 1_G \right) - \left( g_{[1, n]}, g_{[3, n]}, \ldots, g_{[n, n]}, 1_G \right) + \cdots \\
	+ (-1)^n \left( g_{[1, n-1]}, \ldots, g_{[n-1, n-1]} \right) g_n  \\
	= (g_2 | \cdots | g_n)^\dagger + \sum_{k=1}^{n-1} (-1)^k (\cdots | g_k g_{k+1} | \cdots)^\dagger \\
	+ (-1)^n (g_1 | \cdots | g_{n-1})^\dagger g_n.
\end{multline}

% segment-id: o014.li2.chapter6.overview-g-modules-resolutions.p037
The normalized complex is handled in the same way: bases of
$\overline{\mathsf{L}}_n$ as a left or right $G$-module may be taken,
respectively, to be
% segment-id: o014.li2.chapter6.overview-g-modules-resolutions.e007
\begin{equation}\label{eqn:nor-cochain}
	(g_1| \cdots | g_n) \quad \text{or} \quad (g_1 | \cdots | g_n)^\dagger \quad \bmod \mathrm{Triv}_n,
\end{equation}
subject to $g_k\neq1_G$ for every $1\leq k\leq n$.

% segment-id: o014.li2.chapter6.overview-g-modules-resolutions.r002
\begin{remark}
	The resolutions above have a profound topological interpretation:
	$\mathsf{L}$ is the chain complex supplied by the universal torsor
	$\mathrm{E}G$ over the classifying space $\mathrm{B}G$; see
	Examples~\ref{eg:classifying-space-std-cplx} and
	\ref{eg:classifying-contractible}.
\end{remark}
